% Generated draft: requires mathematical and source review.
\documentclass[11pt]{article}
\usepackage[margin=25mm]{geometry}
\usepackage{fontspec}
\setmainfont{Latin Modern Roman}
\usepackage{amsmath,amssymb,mathtools,unicode-math}
\setmathfont{Latin Modern Math}
\usepackage{enumitem,needspace,xurl,hyperref}
\setlength{\parindent}{0pt}
\setlength{\parskip}{5pt}
\setlength{\emergencystretch}{3em}
\newcommand{\N}{\mathbb{N}}
\newcommand{\Z}{\mathbb{Z}}
\newcommand{\Q}{\mathbb{Q}}
\newcommand{\R}{\mathbb{R}}
\newcommand{\C}{\mathbb{C}}
\newcommand{\F}{\mathbb{F}}
\newcommand{\A}{\mathbb{A}}
\newcommand{\PP}{\mathbb{P}}
\newcommand{\E}{\mathbb{E}}
\newcommand{\eps}{\varepsilon}
\newcommand{\dd}{\,\mathrm{d}}
\begin{document}
\section*{UnsolvedMath ID: 1100311}
\textbf{Source problem code:} catalogue entry
\subsection{Definitions and mathematical statement}
Let \(F_n\) be a free group of finite rank \(n\), write \([a,b]=aba^{-1}b^{-1}\), and define
\[
\operatorname{genus}(x)=\min\{g\in\Z_{\ge0}:x=[a_1,b_1]\cdots[a_g,b_g],\ a_i,b_i\in F_n\}
\quad(x\in[F_n,F_n]).
\]
For \(\operatorname{genus}(x)=g\), let
\[
\mathcal E_g(x)=\{(a_1,b_1,\ldots,a_g,b_g)\in F_n^{2g}:\prod_{i=1}^g[a_i,b_i]=x\}.
\]
The requested bound is an explicit finite function \(B(g)\), independent of \(x\), such that
\[
\#\bigl(\mathcal E_g(x)/{\sim_{\mathrm N}}\bigr)\le B(g)
\quad\text{whenever }\operatorname{genus}(x)=g,
\]
where \(\sim_{\mathrm N}\) is the source's Nielsen equivalence of expressions. The excerpt does not define this equivalence. In particular, it does not specify whether one uses arbitrary Nielsen changes of the \(2g\) generators or automorphisms preserving the boundary word \(\prod_{i=1}^g[u_i,v_i]\), nor whether boundary conjugation is allowed. The numerical target cannot be fixed independently of that convention. The statement explicitly requests independence from \(x\); the supplied research summary additionally requests independence from \(n\).

For \(g\ge1\), the supplied partial-progress summary proposes
\[
 f(g)\le W_{\mathrm{tot}}(g)\le\sum_{e=2g}^{6g-3}(2e-1)!!
 \le(4g-2)(12g-7)!!,
\qquad f(g)\le\frac{(12g)!}{(6g)!},
\]
where \((2e-1)!!=\prod_{j=1}^{e}(2j-1)\), and describes \(W_{\mathrm{tot}}(g)\) as counting all Wicks forms. It reports a literature lower bound \(2^g\le f(g)\), describes the displayed upper estimate as roughly \((12g)^{6g}\) up to exponential factors, and distinguishes Bestvina--Feighn's equivalence-class lower bound from Duncan--Vdovina's \(g!\) growth of raw presentations. These are supplied, unverified claims, not conclusions established here. The summary explicitly says its upper bound does not appear in the literature and that its surjectivity input was not checked in Culler's full paper. Its last caveat is truncated. For \(g=0\), the expression set has the unique empty expression for \(x=1\); the displayed positive-genus formula is not applied.
\subsection{Short English statement}
Bound, in terms of the minimal number of commutators alone, the number of Nielsen-equivalence classes of minimal commutator expressions for an element of a free group. The precise equivalence convention and the provenance of the supplied proposed bound require review.
\subsection{Sources}
\begin{itemize}
\item Local UnsolvedMath registry, permanent ID \(1100311\), exact source code \texttt{catalogue entry}, ``Questions in Geometric Group Theory --- Q 3.11.''
\item Bestvina, \emph{Questions in Geometric Group Theory} (2004), Question 3.11, PDF page 12; \url{https://www.math.utah.edu/~bestvina/eprints/questions-updated.pdf}. Manual transcription; recorded access 2026-07-29.
\item Supplied research summary names Culler's 1981 paper, Comerford--Edmunds theory, Bestvina--Feighn, and Duncan--Vdovina, without complete bibliographic citations. It reports checking Culler only through a Crossref record and the role described in Duncan--Vdovina, rather than reading the full article. No authoritative citation for the proposed explicit upper bound is supplied.
\end{itemize}
\end{document}
